%-------------------------- packages ---------------------------%
\usepackage{xcolor}
\usepackage{graphicx}
\usepackage{ifthen}

% %------------------------- font format -------------------------%
% % use helvetica as sans serif font
\usepackage{helvet}
\renewcommand{\familydefault}{\sfdefault}

%--------------------------- colors ----------------------------%

% define JFlyColors
\definecolor{JFlyColor1}{rgb}{0,0,0} % black
\definecolor{JFlyColor2}{rgb}{0.9,0.6,0} % orange
\definecolor{JFlyColor3}{rgb}{0.35,0.7,0.9} % sky blue
\definecolor{JFlyColor4}{rgb}{0.8,0.4,0} % yellowish red
\definecolor{JFlyColor5}{rgb}{0,0.6,0.5} % bluish green
\definecolor{JFlyColor6}{rgb}{0.95,0.9,0.25} % yellow
\definecolor{JFlyColor7}{rgb}{0,0.45,0.7} % blue
\definecolor{JFlyColor8}{rgb}{0.8,0.6,0.7} % reddish purple

%-------------------------- lengthes ---------------------------%
\newlength{\PicWidth}
\newlength{\PicHeight}
\newlength{\ScaleBarEndLength}
\newlength{\LegSymbolLength}
\newlength{\LegSymbolItemSep}
\newlength{\LegItemSep}
\newlength{\tmp}

% scaling
\newcommand{\ScalingFactor}{}

%------------------------- TikZ setup --------------------------%
\usetikzlibrary{calc, shapes} %, decorations.pathreplacing, decorations.pathmorphing, arrows.meta, shapes, backgrounds}
\usepgflibrary{fpu} 
\tikzset{SrfImark/.style={fill=red, draw, isosceles triangle, minimum width=1ex, minimum height=1ex, shape border rotate=-90}}%

%----------------------- Helper commands -----------------------%
\newcommand{\DrawGrid}[1]{
  % x
  \foreach \x in {0,10,...,100}{
    \draw[red, opacity=0.5] ($(#1.north west)!\x/100!(#1.north east)$) node[anchor=south,font={\tiny}] {\x} -- ($(#1.south west)!\x/100!(#1.south east)$);
  }
  \foreach \x in {5,15,...,95}{
    \draw[red, opacity=0.25] ($(#1.north west)!\x/100!(#1.north east)$) -- ($(#1.south west)!\x/100!(#1.south east)$);
  }
  % y
  \foreach \y in {0,10,...,100}{
    \draw[red, opacity=0.5] ($(#1.south west)!\y/100!(#1.north west)$) node[anchor=east,font={\tiny}] {\y} -- ($(#1.south east)!\y/100!(#1.north east)$);
  }
  \foreach \y in {5,15,...,95}{
    \draw[red, opacity=0.25] ($(#1.south west)!\y/100!(#1.north west)$) -- ($(#1.south east)!\y/100!(#1.north east)$);
  }
}

\newcommand{\GetPicWidthHeight}[1]{
    \pgfkeys{/pgf/fpu=true, pgf/fpu/output format=fixed}
    \pgfextractx{\PicWidth}{\pgfpointdiff{\pgfpointanchor{#1}{west}}{\pgfpointanchor{#1}{east}}}
    \pgfextracty{\PicHeight}{\pgfpointdiff{\pgfpointanchor{#1}{south}}{\pgfpointanchor{#1}{north}}}
    \pgfkeys{/pgf/fpu=false}
}

\newcommand{\PrintFigDim}{
    \pgfextractx{\PicWidth}{\pgfpointdiff{\pgfpointanchor{current bounding box}{west}}{\pgfpointanchor{current bounding box}{east}}}
    \pgfextracty{\PicHeight}{\pgfpointdiff{\pgfpointanchor{current bounding box}{south}}{\pgfpointanchor{current bounding box}{north}}}
    \pgfmathsetmacro{\Winches}{\PicWidth/72} % convert to inches
    \pgfmathsetmacro{\Hinches}{\PicHeight/72} % convert to inches
    \pgfmathsetmacro{\Wmm}{\PicWidth/2.835} % convert to mm
    \pgfmathsetmacro{\Hmm}{\PicHeight/2.835} % convert to mm
    \path (current bounding box.south east) node[anchor=south east, red, fill=white] {\pgfmathprintnumber[precision=2]{\Wmm}~mm x \pgfmathprintnumber[precision=2]{\Hmm}~mm};
}